%% ---------------------------------------------------------------------
%%  FIGURA - Cadeia de rastreabilidade requisito-evidencia (DG-09)
%% ---------------------------------------------------------------------
\begin{figure}[htb]
	\caption{Cadeia de rastreabilidade do requisito à evidência}
	\label{fig:rastreabilidade}
	\centering
	\begin{tikzpicture}[
			font=\footnotesize,
			caixa/.style={draw, rounded corners=2pt, minimum height=0.85cm, align=center, fill=black!4, inner sep=4pt},
			>=Stealth,
			fl/.style={->, thick}
		]

		\node[caixa] (mem) at (0,0) {Solicitação\\ do produto};
		\node[caixa, right=0.6cm of mem] (req) {Requisitos\\ \texttt{FR}/\texttt{NFR}};
		\node[caixa, right=0.6cm of req] (ca) {Critérios\\ \texttt{CA}};
		\node[caixa, below=0.9cm of ca] (adr) {Decisões\\ \texttt{ADR}};
		\node[caixa, left=0.6cm of adr] (tar) {Tarefas\\ \texttt{T}};
		\node[caixa, left=0.6cm of tar] (cod) {Código\\ artefato};
		\node[caixa, below=0.9cm of cod] (evd) {Evidência\\ execução e inspeção};
		\node[caixa, right=0.6cm of evd] (res) {Resultado\\ PASS / PENDENTE};

		\draw[fl] (mem) -- (req);
		\draw[fl] (req) -- (ca);
		\draw[fl] (ca) -- (adr);
		\draw[fl] (adr) -- (tar);
		\draw[fl] (tar) -- (cod);
		\draw[fl] (cod) -- (evd);
		\draw[fl] (evd) -- (res);
	\end{tikzpicture}
	\legend{Fonte: elaborado pelo autor (2026). A ausência de critério rompe a cadeia entre requisito e
	evidência: é exatamente o caso dos dez requisitos marcados como \textbf{SEM CA} no Apêndice~\ref{ap:rastreabilidade}.}
\end{figure}
